% !TEX root=schubert.tex
\documentclass[document.tex]{subfiles}
\usepackage{xcolor}

\begin{document}
\section{\cite{scto01} basic equations}
This document/section has a different approach to structure the derivation of the ALA: this is a summary of the derivation of the ALA done by \cite{scto01}.
\textcolor{red}{red text means assumtions/approximations were made toward the ALA}. 
\subsection{Mass equation/continuity equation/conservation of mass}
\begin{align} \label{mass scto}
\frac{\partial \rho}{\partial t} + \frac{\partial}{\partial x_1}(\rho u_i) = \frac{1}{\rho}\frac{D\rho}{Dt} + \frac{\partial u_1}{\partial x_1} = 0
\end{align}

where $\frac{D}{Dt}$ represents the material derivative (\ref{material derivative}). For incompressible flow, $\frac{D\rho}{Dt} = 0$, so eq \ref{mass scto} reduces to: 
\begin{align} \label{mass equation}
\frac{\partial u_i}{\partial x_i} = \nabla \cdot \textbf{u} =0
\end{align}

\textcolor{red}{For the special case of an incompressible fluid, $\mdv{\rho}{t} = 0$, so equation \ref{mass equation} reduces to}
\begin{align}
    \color{red}
    \pdv{u_i}{x_i} = 0
\end{align}

\subsection{Momentum equations/conservation of momentum}
Starting governing equation (6.4.1): 
\begin{align} \label{momentum equation}
    \rho \mdv{u_i}{t} = - \pdv{p}{x_i} + \pdv{\tau_{ij}}{x_j}+\rho g_i
\end{align}
\textcolor{red}{we assume gravity $g$ to be the only body force.The coriolis force is negligibly small in the Earth's mantle, and the centrifugal 
force due to the Earth's rotation is included in the gravity term.}

In 2D, the momementum equation components are: 
\begin{align}
\rho \left(
    \pdv{u_x}{t}
    + u_x \pdv{u_x}{x}
    + u_y \pdv{u_x}{y}
\right)
&= -\pdv{p}{x}
   + \pdv{\tau_{xx}}{x}
   + \pdv{\tau_{xy}}{y},
\\
\rho \left(
    \pdv{u_y}{t}
    + u_x \pdv{u_y}{x}
    + u_y \pdv{u_y}{y}
\right)
&= -\pdv{p}{y}
   + \pdv{\tau_{xy}}{x}
   + \pdv{\tau_{yy}}{y}
   + \rho g.
\end{align}

\textcolor{red}{Note that the assumption is made that gravity acts in the positive y-direction. Section 6.14?ext}

\subsection{Navier-Stokes equation}
The local strain rate tensor $e_{ij}$ is given by: 
\begin{align} \label{strain rate tensor}
    e_{ij} = \frac{1}{2}\left(\pdv{u_i}{x_j}+\pdv{u_j}{x_i}\right)
\end{align}
\textcolor{red}{For a Newtonian, isotropic fluid, the relation between the deviatoric(!) stress tensor and the strain rate tensor is}
\begin{align} \label{6.5.2}
    \color{red}
    \tau_{ij} = 2\mu e_{ij} + \lambda e_{kk} \delta_{ij}
\end{align}

Defining the bulk viscosity $k_B$: 
\begin{align} \label{6.5.3}
    \frac{\tau_{ii}}{3} = e_{ii}\left(\lambda +\frac{2}{3}\mu\right) = k_Be_{ii}
\end{align}
The combination of equations \ref{6.5.2}, \ref{6.5.3} and \ref{strain rate tensor} gives
\begin{align} \label{stress tensor with bulk viscosity}
    \tau_{ij} = \mu \left(\pdv{u_i}{x_j}+\pdv{u_j}{x_i}\right) + \left(k_B - \frac{2}{3} \mu \right)\pdv{u_k}{x_k}\delta_{ij}
\end{align}

and substitution of eq. \ref{stress tensor with bulk viscosity} into eq. \ref{momentum equation} gives
\begin{align} \label{6.5.5}
    \rho \mdv{u_i}{t} = - \pdv{p}{x_i} + \pdv{x_j}\left[\mu \left(\pdv{u_i}{x_j}+\pdv{u_j}{x_i}\right) + 
    \left(k_B - \frac{2}{3} \mu \right)\pdv{u_k}{x_k}\delta_{ij}\right] + \rho g_i
\end{align}

\textcolor{red}{For many fluids, $k_B$ is taken to be zero \textbf{(Stokes assumption)}, 
so eq. \ref{stress tensor with bulk viscosity} becomes} 

\begin{align} \label{stokes stress tensor}
    \color{red}
     \tau_{ij} = \mu \left(\pdv{u_i}{x_j}+\pdv{u_j}{x_i} - \frac{2}{3}\pdv{u_k}{x_k}\delta_{ij}\right)
\end{align}

\textcolor{red}{and eq. \ref{6.5.5} becomes the Navier-Stokes equation:}
\begin{align} \label{6.5.7}
    \color{red}
    \rho \mdv{u_i}{t} = - \pdv{p}{x_i} + \pdv{x_j}\left[\mu \left(\pdv{u_i}{x_j}
    +\pdv{u_j}{x_i}-\frac{2}{3}\pdv{u_k}{x_k}\delta_{ij}\right)\right]
    + \rho g_i
\end{align}


\textcolor{red}{for incompressible flow, equation \ref{mass equation} tells us $\pdv{u_k}{x_k} =0$, and together with a constant dynamic viscosity $\mu$ so eq. \ref{6.5.7} reduces to}: 
\begin{align}
    \color{red}
    \rho \mdv{u_i}{t} = - \pdv{p}{x_i} + \mu \pdv[2]{u_i}{x_j} + \rho g_i
\end{align}

This equation represents a balance between inertial forces $\left(\mdv{u_i}{t}\right)$, pressure forces $\left(\pdv{p}{x_i}\right)$, viscous forces $\left(\mu \pdv[2]{u_i}{x_j}\right)$ 
and the body force due to gravity $\left(\rho g_i\right)$

\subsection{Thermal energy equations/conservation of energy}
\cite{scto01} presents various different formulations of the energy equation.

The general form of the energy equation (6.9.1) is:
\begin{equation} \label{6.9.1}
    \rho T \mdv{s}{t} = \tau_{ij}\pdv{u_i}{x_j} + \pdv{}{x_j}\left(k\pdv{T}{x_j}\right) + \rho H 
\end{equation}
The terms on the right side of \ref{6.9.1} are, respectively, volumetric heat production rates due to viscous dissipation (internal friction), thermal conduction,
and internal heat generation (radiogenic?). 
\textcolor{red}{Where it is assumed that Fourier's law of heat conduction for an isotropic medium is applicable:}
\begin{align}
    \color{red}
    q_i = -k \pdv{T}{x_i}
\end{align}
Where $q_i$ is the heat flux vector. Additionally, the viscous dissipation term can be written as the dissipation function: 
\begin{align} \label{dissipation function}
    \Phi = \tau_{ij} \pdv{u_i}{x_j}
\end{align}
Inserting the definition of our deviatoric stress tensor with $k_B = 0$ (eq. \ref{stokes stress tensor}), this becomes
\begin{align} \label{6.9.4}
    \Phi = 2\mu \left\{\frac{1}{2}\left(\pdv{u_i}{x_j} + \pdv{u_j}{x_i}\right)
    - \frac{1}{3}\left(\pdv{u_k}{x_k}\right)\delta_{ij}\right\}^2
\end{align}

\textcolor{red}{For an incompressible and Newtonian fluid ($\pdv{u_i}{x_i} = 0$), the dissipation function reduces to}
{\color{red}
\begin{align}
    \Phi = \frac{\mu}{2}\left(\pdv{u_i}{x_j} + \pdv{u_j}{x_i}\right)^2
\end{align}
}

With various thermodynamic relationships: 
\begin{align}
    \alpha &= \frac{1}{v}\left(\pdv{v}{T}\right)_p \\
    c_p &= T\left(\pdv{s}{T}\right)_p \\
    \left(\pdv{s}{p}\right)_T &= -\left(\pdv{v}{T}\right)_p
\end{align}
It can be shown that 
\begin{align} \label{6.9.6}
    \mdv{s}{t} = \frac{c_p}{T}\mdv{T}{t} - \frac{\alpha}{\rho} \mdv{p}{t}
\end{align}

Plugging \ref{6.9.6} into \ref{6.9.1} gives 
\begin{align} \label{6.9.7}
    \rho c_p \mdv{T}{t} - \alpha T \mdv{P}{t} &= \pdv{x_i} \left(k\pdv{T}{x_i}\right) + \Phi + \rho H \\
    \rho c_p \left(mdv{T}{t} - \frac{\alpha T}{\rho c_p} \mdv{P}{t} \right) &=\pdv{x_i} \left(k\pdv{T}{x_i}\right) + \Phi + \rho H \\
\end{align}

% OLD CODE %
% For an incompressible Newtonian fluid (eq. 6.9.13 + 6.9.14): 
% \begin{equation}
%     \rho C_v \mdv{T}{t} = \pdv{}{x_i}(k\pdv{T}{x_i}) + \frac{\mu}{2}\left(\pdv{u_i}{x_j} + \pdv{u_j}{x_i}\right)^2+\rho H
% \end{equation}
% with $\frac{\mu}{2}\left(\pdv{u_i}{x_j} + \pdv{u_j}{x_i}\right)^2$ the viscous dissipation function. 

% In terms of $C_p$ instead of $C_v$ (for a compressible fluid) (eq. 6.9.16): 
% \begin{equation}
%     \rho C_p \left(\mdv{T}{t} - (\pdv{T}{p})_s \mdv{p}{t}\right) = \pdv{}{x_i}(k\pdv{T}{x_i}) + \frac{\mu}{2}\left(\pdv{u_i}{x_j} + \pdv{u_j}{x_i}\right)^2+\rho H
% \end{equation}
% with (eq. 6.9.15, 6.9.17):
% \begin{align*}
%     (\pdv{T}{p})_s &= \frac{\alpha T}{\rho C_p} \\
%     (\mdv{T}{t})_s &= (\pdv{T}{p})_s \mdv{p}{t} = \frac{\alpha T}{\rho C_p} \mdv{p}{t}
% \end{align*}

% \begin{align}
% \rho \mdv{u_i}{t} = -\pdv{p}{x_i} + \pdv{\tau	_{ij}}{x_j} + \rho g_i
% \end{align}

\section{Approximate equations}





\end{document}